The functions associated with the poly-mod section are very powerful. They provide the means for non-linear modulation: cross-modulattion of the two voice oscillators and also to modulation the filter cut-off frequency by VCO B. 

The familiar and common synthesizer feature is that the oscillator (frequency,filter  cut-off frequency, pulse width, etc.) is modulated by a \textit{low frequency oscillator (LFO)}. An LFO does not operate in the audible frequency range (it is "low frequency") and it has a fixe frequency, e.g. independent from the note played. Poly-mod is different in that the modulation itself is in the audible range and the resulting non-linearity changes the frequency spectrum drastically. The resulting effect is compareable to FM synthesis known from digital synthesizers, although in the Prophet 600 it is more raw and fully analog. Like the oscillator sync function, cross-modulation is hardware implemented.

The poly-mod control can be found in the poly-mod sub-panel. 

\begin{center}
\scalebox{0.4}{
  \begin{tikzpicture}[scale=0.8]

    \addpar{-18cm,3.5cm}{\envrouting \\ The routing options determine whether poly-mod is driven by the filter or the 2nd envelope};
    \draw[line width = 2pt, dashed] (-4cm,3.5cm) -- (1cm,3.5cm);


   \polymodesubpanel{0,0}
  \end{tikzpicture}
}
\end{center}

There are two modulation sources, \polyenv and \oscamt and two destinations, oscillator a pitch and filter cut-off frequency, which are activated using the respective panel switches \plyfreqa and \plyfilter. The \oscamt control determines how strongly the oscillator B modulates poly-mod targets (destinations). It is instructive to consider the effect of oscillator B modulation on the two targets first and then explore the effect of envelope modulation.

\textbf{Cross-modulation}

\begin{itemize}
  \item When setting the \plyfreqa switch to \textit{on}, oscillator B modulates the frequency of oscillator A, e.g. frequency modulation (FM). Depending on the modulation strength the result is a wide and rich frequency spectrum which contains among others the sum and difference of multiples of the two oscillators frequencies . The spetrum strongly depends on the frequency difference and may be totally non-harmonic but also harmonic for carefully adjusted frequency settings. Typical spectra are often described as bell sounding. Note that in contrast to FM synthesizers which use highly accurate sin waves, the oscillator output of the Prophet does not have the same basic accuracy and the available waveforms already contain higher harmonics. Therefore the FM result in the Prophet 600 is typically more drastic and dirtier and requires classic subtractive methods for shaping the sound, e.g. the application of the low pass filter. 
  \item When setting the \plyfilter switch to \textit{on}, oscillator B modulates the filter cut-off frequency of the VCF. Like frequency modulation this is a non-linear modulation. This effect can also produce a wide frequency spectrum which may or may not be harmonic, depending on the relative frequency of oscillators A and B. The effect is more pronunced at low filter cutt-off frequencies. The resonance also has a strong effect.  
\end{itemize}

Note that poly-mod is very tuning sensitive and that it maybe necessary to retune the synth more often when using extreme poly-mod settings.

\textbf{Envelope pitch modulation}

The second modulation source in the poly-mod envelope whose strength set by the \polyenv dial. This activates a pitch modulation of oscillator A using an envelope when the Frequency switch is \textit{on}. The \polyenv control is bipolar, that is it provides negative (inverse modulation) and positive values. Note that the poly-mod envelope only affects the frequency of oscillator A. The filter cut-off frequency is already modulated by the filter envelope. 

The pitch modulation by envelope is not a cross-modulation function. However, the association with poly-mod is useful because of the drastic effect of frequency (differences) on the sound spectrum. By time modulating the pitch, the sound can be made to sweep through different spectra, ranging from apocalyptic noise through bell sounds to delicate transients. The way humans perceive natural sounds is strongly influenced by the frequency characteristics of the transients (typically short lived high frequency components) in the attack and decay phases of the sound. This is best modelled using an envelope making it a versatile tool for advanced sound design. 

IN order to further expand the sound design posisbilities, the latest firmware upgrade offers envelope routing   option. The menu parameter \envrouting determines if the filter envelope or the assignable envelope is used in poly-mod\footnote{Up to version 2.1 RC3 the filter envelope was "hard wired" to poly-mod.}. The following figure provides an overview of the routing options. Having different envelopes modulating the filter and the poly-mod opens up new possibilites in the sound and spectrum dynamics. 

\begin{itemize}
  \setlength\itemsep{0cm}
  \item \textit{standard}: In this "standard" routing the amplitude is modulated by the 2nd envelope and the poly-mod is modulated by the filter envelope. This is the setup of the original Prophet 600. 
  \item \textit{poly-amplitude}: In this routing both the amplitude and the poly-mod are modulated by the 2nd envelope. 
  \item \textit{poly}: In this routing the amplitude is modulated by the filter envelope and the poly-mod is modulated exclusively by the 2nd envelope. Therefore the poly-mod time modulation is completely free, at the expense of having filter modulation and amplitude tied together.
  \item \textit{gate}: In this routing the amplitude is not modulated at all, but behaves like a gate (similar to electronic organs). The poly-mod is modulated exclusively by the assignable envelope. Therefore the poly-mod modulation is completely free at the expense of have a less dynamic sound evolution.
\end{itemize}  

\scalebox{0.4}{
  \begin{tikzpicture}[scale=0.8]
   \envroute{0,0}
  \end{tikzpicture}
}

\textbf{Oscillator sync and poly-mod}

As explained in section \ref{osc}, oscillator A can be sync'ed to oscillator B using the \oscsync switch in the sub-panel for oscillator B. This can also be view as a variant of cross-modulation. The sync function can even be combined with frequency modulation as decribed above. 

Combining oscillator sync with poly-mod pitch modulation by envelope can be used to produce interesting short transients or slow sound evolutions. For example, when oscillator A is a pulse wave, then sync'ing this wave to oscillator B and time modulating the pitch of oscillator A has the effect of pulse width modulation by envelope. It creates a change in timbre which can be short to simulate plucking an brassy attacks or long to produce a sweep of the sound spectrum.

Note: the firmware versions 2.0 and 2.1 RC3 contain a software error which incorrectly reset the pulse widths of oscilators A and B. The manifests itself under the following conditions: \oscamt is turned up and the poly-mod target \plyfreqa is active. The wave shape pulse of oscillator B is not activated. In this case the incorrect reset of the pulse widths leads to a wrong frequency in oscillator A. The effect becomes strong with higher settings of \oscamt. In version \version this bug was corrected but in order to ensure backwards compatibility of older patches, the bug can be switched off in the menu patch parameters: the parameter \syncbug can be set to \textit{off} and \textit{on}. The default value is \textit{off} but for patches from older firmware versions it is set to \textit{on} so that these are unchanged even if the conditions above apply.    

\textbf{Starting points for patch design with poly-mod and sync}

Using poly-mod in combination with oscillator sync, different modulations strengths for different oscillator pitches with different filter and modulation shapes is a ver wide field of experimentation. Here are some ideas as potential starting points

\begin{itemize}
  \item Set the envelope routing to \textit{standard}, tune both oscillators to the same frequency, in the poly-mod scetion turn source VCO B to zero and deactivate the filter modulation. Then set VCO A to sync and experiment with the poly-mod envelope source amount. This produces a sync whic, for high envelope amounts, can be tuned to a quite harsh synth solo sound.
  \item ... (to be done)
\end{itemize} 
